\section{The abstract structure theorem (Layer~1)}\label{sec:programme}

The bridge (\Cref{sec:bridge}) produces an $\eps$-JB algebra; the factorization
routes (\Cref{sec:factorization,sec:classical}) bypass the abstract structure
theory for the channel application. This section concerns the standalone
\emph{abstract} structure theorem --- Kitaev's main theorem, Jordan-ised --- which
says that an $\eps$-JB algebra is, up to a controlled deformation, a genuine
JB-algebra. It is open. We state the target, the cohomological strategy, the one
quantitative obstruction, and the substantial recent progress: a dimension-free
\emph{exact-adjoint splitting benchmark} that resolves the obstruction at the
level of exact cocycles for every simple factor, together with a precise account
of what still separates the benchmark from the full theorem.

\subsection{The target}

% PROV: Kitaev2024 approximate_algebras.tex:460-462 (th_main) -- C*-statement; ours is the declared Jordan analogue (open).
\begin{openproblem}[Jordan structure theorem]\label{op:jordan-structure}
\status{(open)}
There are universal constants $C,\eps_{0}>0$ \emph{(independent of dimension)}
such that for every $\eps<\eps_{0}$ and every finite-dimensional $\eps$-JB algebra
$\mathcal{A}$ there exist a genuine finite-dimensional JB-algebra $\mathcal{B}$
and a linear bijection $v\colon\mathcal{B}\to\mathcal{A}$ that is a
$C\eps$-Jordan-isomorphism:
\[
   \norm{v(\Id_{\mathcal{B}})-\Id_{\mathcal{A}}}\le C\eps,\qquad
   \norm{v(b\bullet b')-v(b)\bullet v(b')}\le C\eps\,\norm{b}\,\norm{b'},
\]
and $(1-C\eps)\norm{b}\le\norm{v(b)}\le(1+C\eps)\norm{b}$ for all
$b,b'\in\mathcal{B}$.
\end{openproblem}

This is Kitaev's statement~\cite[Theorem \texttt{th\_main}]{Kitaev2024} with the
associative product replaced by the Jordan product and ``$C^{*}$-algebra'' by
``JB-algebra''. What we can state is the strategy and the precise point where the
dimension-free constant is in question.

\subsection{Strategy and vocabulary}

Kitaev builds the model algebra $\mathcal{B}$ incrementally and controls the
error at each step cohomologically. The Jordan programme mirrors four parts.
An \emph{idempotent element} is $e$ with $e\bullet e=e$; a \emph{Jordan frame} is
a list of nonzero pairwise-orthogonal idempotents summing to $\Id$ (the analogue
of diagonal matrix units); the \emph{Peirce decomposition} of an idempotent $e$
splits $\mathcal A$ into the $1,\tfrac12,0$ eigenspaces of $L_e\colon a\mapsto
e\bullet a$; a \emph{derivation} is a linear $D$ with
$D(a\bullet b)=Da\bullet b+a\bullet Db$ (the $D$ form the Lie algebra of
$\Aut(\mathcal A)$).
\begin{enumerate}[label=\textbf{(\alph*)},leftmargin=*]
  \item \textbf{Approximate idempotents / Jordan frame.} Produce elements with
    $\norm{e\bullet e-e}$ small, not near $0$ or $\Id$. The spectral idempotent
    $P$ of \Cref{sec:spectral} is a \emph{map-level} device; it does not by
    itself produce idempotent \emph{elements} of $A$, so this is a separate open
    step (in the channel application one instead uses the ambient $\Bsa$).
  \item \textbf{Approximate Peirce decomposition.} For an approximate idempotent
    $L_e$ has spectrum clustered near $\{0,\tfrac12,1\}$, giving an approximate
    Peirce splitting with $O(\eps)$ cross-terms.
  \item \textbf{Approximate coordinatization.} On a frame of rank $\ge3$ the
    Jordan--von Neumann--Wigner coordinatization~\cite{JordanVonNeumannWigner1934}
    identifies the algebra with Hermitian matrices over a coordinate $*$-algebra;
    the approximate version reads off coordinates up to $O(\eps)$, producing a
    candidate $\mathcal B$ and a linear inclusion $v\colon\mathcal B\to\mathcal A$
    multiplicative up to a (possibly large) $\delta$.
  \item \textbf{Error reduction.} Replace $v$ by a nearby map whose defect is
    $O(\eps)$, \emph{uniformly in dimension}. This is the cohomological heart,
    treated below.
\end{enumerate}

\subsection{The cohomological error reduction}

% PROV: ORIGINAL agent-A/theory/01-error-reduction.md; agent-B/notes/next-arrow-to-newton-error-reduction.md -- Newton step, splitting, next-arrow.
Write $g(b,b')=v(b\bullet b')-v(b)\bullet v(b')$ for the multiplicativity defect
of the inclusion $v$ from step (c); it is symmetric with $\norm{g}\le\delta$.
Linearizing the Jordan identity ($\mathcal B$ exact, $\mathcal A$ only $\eps$-JB)
shows $g$ is an \emph{$\eps$-approximate Jordan $2$-cocycle}. The Jordan
$2$-coboundary of a linear $h\colon\mathcal B\to\mathcal A$ is
\[
  (d^{1}h)(b,b')=v(b)\bullet h(b')+h(b)\bullet v(b')-h(b\bullet b').
\]
All cochain norms below are injective order-unit norms:
\[
  \norm{\theta}_{\mathrm{inj}}
   =\sup_{\norm{x_i}\le1}\norm{\theta(x_1,\dots,x_k)} .
\]
We write $\operatorname{dist}(\theta,\Img d^1)$ for the distance in this norm.
For an exact product $\jp$ and a unit-normalized symmetric $2$-cochain
$\theta(\Id,x)=0$, the linearized Jordan-identity defect is
\[
\begin{aligned}
  (J\theta)(a,b)
  ={}&\theta(a^2,a\jp b)+a^2\jp\theta(a,b)
      +(a\jp b)\jp\theta(a,a) \\
    &-\theta(a^2\jp b,a)-a\jp\theta(a^2,b)
      -a\jp\bigl(b\jp\theta(a,a)\bigr).
\end{aligned}
\]
This is the first-order part of
$(a*a)*(b*a)-((a*a)*b)*a$ for $x*y=x\jp y+\theta(x,y)$; the informal symbol
$d^2\theta$ in older internal notes is not used as a separate cohomological
differential here. Below the actual first-order operator is always $J\theta$
(with quadratic error terms handled separately in the Newton step).
Two ingredients drive a Newton iteration:
\begin{itemize}[leftmargin=*]
  \item a \emph{splitting} (right inverse) $s$ of $d^1$ with
    $\norm{s}\le K$: a recipe $h=sg$ with $d^1h=g$ on exact cocycles;
  \item a \emph{next-arrow (cocycle-projection) estimate}: every approximate
    cocycle $\theta$ is within $K'\norm{J\theta}$ of an exact one, so that
    $g$ may be corrected.
\end{itemize}
Setting $\tilde v=v-h$ replaces $\delta$ by
$O(\delta^2+(K+K')\eps)$; iterating the Newton step
\[
  \delta\ \longmapsto\ C(K,K')\,(\eps+\delta^2)
\]
drives small $\delta$ to $O(\eps)$. \emph{The theorem reduces to: are $K,K'$
dimension-free, in the order-unit norm?} (The Newton bookkeeping itself is
verified, conditional on $K,K'$.)

The two qualitative algebraic inputs are classical.

% PROV: ChuRusso2016 (Thm 1.1/1.2) via agent-A/notes/ingestion-results-2026-06-01.json; Penico1951. Extraction-level provenance; primary PDF byte-check pending.
\begin{theorem}[Whitehead lemmas for Jordan algebras]\label{thm:whitehead}
\status{(cited; source-extracted, primary byte-check pending)}
Let $J$ be a finite-dimensional semisimple Jordan algebra over a field of
characteristic $0$ and $M$ a finite-dimensional $J$-module. Then every derivation
$J\to M$ is inner ($H^{1}(J,M)=0$, Jacobson); and if $J$ is separable, every
Jordan $2$-cocycle is a coboundary ($H^{2}(J,M)=0$, Penico).
\end{theorem}

% PROV: FarautKoranyi1994 via agent-A/notes/ingestion-results-2026-06-01.json (item 5). Extraction-level provenance.
\begin{proposition}[Compactness of $\Aut(J)$]\label{prop:aut-compact}
\status{(cited; source-extracted, primary byte-check pending)}
For a finite-dimensional Euclidean Jordan algebra $J$, $\Aut(J)$ is a compact
subgroup of the orthogonal group of the trace inner product; every automorphism
preserves spectra, hence the order-unit norm. Normalized Haar averaging over
$\Aut(J)$ is therefore available.
\end{proposition}

The vanishing $H^{1}=H^{2}=0$ makes the correction step possible \emph{in
principle}; the issue is its \emph{size}. Crucially, Haar averaging over
$\Aut(J)$ is a norm-one \emph{projection} onto invariants but is \emph{not} a
right inverse to $d^1$; the inverse on the non-invariant cochain components is
exactly where dimension can enter.

\subsection{The rank obstruction}

The single quantitative difficulty can be named precisely.

% PROV: ORIGINAL agent-A/theory/02-spin-splitting.md (rank obstruction), agent-B/notes/cochain-norm-conversion-caveat.md.
\begin{proposition}[The order-vs-Frobenius gap is exactly $\sqrt{\mathrm{rank}}$]
\label{prop:rank-gap}
\status{(proved)}
On a simple Euclidean Jordan algebra of rank $r$, with order-unit norm
$\norm{a}_{\mathrm{op}}=\max_i\lvert\lambda_i(a)\rvert$ and Frobenius norm
$\norm{a}_2=(\sum_i\lambda_i(a)^2)^{1/2}$,
\[
  \norm{a}_{\mathrm{op}}\le\norm{a}_2\le\sqrt{r}\,\norm{a}_{\mathrm{op}},
\]
and the upper ratio $\sqrt r$ is attained at $a=\Id$. Thus
$r=2$ for spin factors (dimension-free), $r=n$ for $H_n(\R/\C/\mathbb H)$
(growing), $r=3$ for the Albert algebra.
\end{proposition}

The cohomological splitting is known to be dimension-free in the \emph{Frobenius}
norm: the Casimir idempotent of the multiplication algebra $M(J)$, written as a
Haar average over the compact unitary group $U(M(J))$ of
\emph{Frobenius}-isometries, has projective tensor norm $\le1$, giving
$\norm{s}_{\mathrm{F}}=1/\sigma_{\min}(d^1)=O(1)$ (confirmed numerically:
$\dim\Ker d^1=\dim\mathrm{Der}(J)$ matched exactly, and $1/\sigma_{\min}$ stays
$\le1$ and decreasing for $H_n(\R/\C)$ and spin out to large $n$). This is
\emph{not} enough: $U(M(J))$-isometries are not order-isometries
(\Cref{prop:aut-compact} gives only $\Aut(J)$ for the order norm), so the
Frobenius bound need not survive the conversion to the order-unit norm, with a
worst-case loss of $\sqrt r$. Two caveats forbid cheap fixes:

% PROV: agent-B/notes/cochain-norm-conversion-caveat.md, nuclear-rank-one-route-caveat.md.
\begin{remark}[Why a Frobenius splitting cannot simply be converted]\label{rem:cochain-caveats}
\status{(proved; obstruction)}
On $H_n(\R)$ the linear map $h(x)=x_{11}\Id$ has $\norm{h}_{\mathrm{op}}=1$,
$\norm{h}_{\mathrm F}=\sqrt n$, and an order-bounded coboundary
$\norm{d^1h}_{\mathrm{op}}\le3$, yet every primitive of $d^1h$ has Frobenius norm
$\ge\sqrt n$. So an order-bounded coboundary can force a $\sqrt n$ Frobenius
primitive: one must not seek Frobenius-bounded primitives, only order-bounded
ones. Likewise the rank-one decomposition route fails: the identity on the
trace-zero part of $H_n(\C)$ has nuclear norm $\ge\dim\ge n^2$ while its operator
norm is $1$, so summing rank-one primitive estimates loses dimension. The
splitting must be estimated \emph{directly} in the order norm.
\end{remark}

\subsection{The exact-adjoint splitting benchmark}

The following results (Agent~B, 2026-06-05) resolve the obstruction at the level
of \emph{exact} cocycles in the \emph{adjoint} module $M=\mathcal B$: a
dimension-free order-unit-norm splitting of $d^1$ exists for every simple factor
and, by a summand-count-free direct-sum reduction, for every finite-dimensional
JB-algebra. This is a major step; it is \emph{not} yet \Cref{op:jordan-structure}
(see \Cref{ssec:layer1-open}). We mark these
\status{(proved; exact-adjoint benchmark)}.

% PROV: agent-B/notes/adjoint-spin-splitting-theorem.md, agent-A/theory/02-spin-splitting.md (independent confirmation).
\begin{proposition}[Spin factors]\label{prop:spin-splitting}
\status{(proved; exact-adjoint benchmark)}
For a spin factor $V=\R\Id\oplus H$ there is an $O(H)$-equivariant right inverse
$S$ of $d^1$ (adjoint module) with
$\norm{Sf}_{\mathrm F}\le2\norm{f}_{\mathrm{inj}}$;
by the rank-$2$ norm comparison (\Cref{prop:rank-gap}) the order-unit constant is
$\norm{S}_{\mathrm{op}\to\mathrm{op}}\le4\sqrt2$, \emph{independent of}
$\dim H$. The associated next-arrow estimate also holds: a normalized adjoint
$2$-cochain $\theta$ with $\theta(\Id,z)=0$ obeys
$\operatorname{dist}(\theta,\Img d^1)\le(2\sqrt2+2)\norm{J\theta}$.
\end{proposition}

\noindent (Independently obtained by both authors; the rank-$2$ structure makes
the order and Frobenius norms $\sqrt2$-equivalent, so the spin gap of
\Cref{prop:rank-gap} is absent.)

% PROV: agent-B/notes/adjoint-direct-sum-reduction.md, spin-direct-sum-adjoint-corollary.md, bounded-rank-adjoint-reduction.md; agent-B/notes/response-to-agent-a-v0.12-layer1-caveat.md (mixed Peirce-1/2 caveat).
\begin{proposition}[Direct sums]\label{prop:direct-sum}
\status{(proved; exact-adjoint benchmark)}\afbadge{prop-direct-sum}
If $\mathcal B=\bigoplus_{r}\mathcal B_r$ (order-unit norm $=\max_r$) and each
factor has an exact adjoint splitting with constant $K_r$, then $\mathcal B$ has
one with constant $\max_r K_r+1$, \emph{independent of the number of summands}
(off-block components are recovered by $P_r f(e_r,\cdot)$, with no sum over $r$).
\end{proposition}

% PROV: agent-B/notes/response-to-agent-a-v0.12-layer1-caveat.md.
\begin{remark}[Scope of summand-count independence]\label{rem:directsum-scope}
\status{(proved)}
\Cref{prop:direct-sum} is for the adjoint (or block-respecting) module. For an
\emph{arbitrary} Jordan module the coboundary need not be block-diagonal even
when cross products vanish in $\mathcal B$: central idempotents can act with
Peirce eigenvalue $\tfrac12$ on a module component (e.g.\ for $\R e_1\oplus\R e_2$
the off-diagonal Peirce module of $H_2(\R)$ has both $e_i$ acting by $\tfrac12$).
Such mixed Peirce-$\tfrac12$ components must be controlled separately; this is
the module-side qualification on summand-count independence.
\end{remark}

% PROV: agent-B/notes/commutative-scalar-module-splitting.md, commutative-scalar-cocycle-projection-theorem.md.
\begin{proposition}[Commutative scalar modules]\label{prop:comm-scalar}
\status{(proved)}
Algebraically, every finite-dimensional unital Jordan $\R^m$-module splits into
one-dimensional modules with action $x\cdot m=\ell(x)m$, where $\ell(x)=x_k$
(coordinate) or $\ell(x)=\tfrac12(x_p+x_q)$ (half-sum). In the associated
max-sector norm, the support-unit formula $Sf(x)=f(x,s)$ (with $s=e_k$ resp.\
$e_p+e_q$) is a norm-one right inverse, and $\Pi=d^1S$ is a projection onto
$\Img d^1$ with $\norm{\Pi}\le3$. For arbitrary equivalent module norms this
statement carries the corresponding complementability constant. In the
max-sector norm the next-arrow estimate
$\operatorname{dist}(\theta,\Img d^1)\le12\norm{J\theta}$ holds (sharpened to
$\le2$ for coordinate modules), including the Peirce-$\tfrac12$ half-sum
modules.
\end{proposition}

The hard case is the high-rank matrix factors, where the rank gap of
\Cref{prop:rank-gap} is real.

% PROV: agent-B/notes/matrix-factor-exact-adjoint-splitting-theorem.md, fixed-frame-peirce-matrix-reduction.md, off-sector-leakage-globalization-theorem.md, real-symmetric/pure-skew/complex-antilinear/quaternionic reconstruction theorems, frame-covariance-and-global-matrix-obstacle.md.
\begin{theorem}[High-rank matrix factors]\label{thm:matrix-splitting}
\status{(benchmark; Agent B proof, re-audit pending)}
There is a universal constant $C$ such that for every $n$ and
$F\in\{\R,\C,\mathbb H\}$: for an exact adjoint coboundary $f=d^1h$ on
$J=H_n(F)$ there is a derivation $\delta\in\mathrm{Der}(J)$ with
$\norm{h-\delta}\le C\norm{f}_{\mathrm{inj}}$, i.e.\ $d^1$ has an order-unit
right inverse of norm $\le C$ \emph{uniformly in $n$}. The construction proceeds
by: a fixed diagonal-frame Peirce gauge ($\norm{h|_D}\le11\norm{f}$); a coherent
off-sector leakage globalization (averaged squared diagonal-sign commutator with
spectral gaps $\{0,\tfrac12,1,\tfrac32\}$); and sector-preserving edge-map
reconstruction in each Peirce sector --- real-symmetric (constant $60$), complex
Hermitian ($84$), complex antilinear ($124$), and quaternionic --- via
random-triple matching-curvature averaging. The fixed-frame splitting is
frame-covariant (constant $\le11$, projection $\le33$ for every Jordan frame).
\end{theorem}

% PROV: agent-B/notes/subagent-schur-residual-audit-v0.1.md; subagent-matching-curvature-audit-v0.1.md; pointwise-schur-curvature-caveat.md.
\begin{remark}[Verification status of \Cref{thm:matrix-splitting}]\label{rem:matrix-audit}
\status{(audit caveat)}
The matrix construction is recent and intricate. Two sidecar audits examined
\emph{earlier} formulations of the sector-reconstruction step (a
Haagerup/base-vertex route) and found a genuine ordinary-vs-completely-bounded
gap there; Agent~B's final proof instead uses random-triple matching
reconstruction, which is designed to avoid that gap, but the audits did not
re-verify the final argument. A logarithmic Schur-multiplier family is the
sharpest stress test; its coboundary appears to grow logarithmically as well, so
it is not a counterexample. We therefore record \Cref{thm:matrix-splitting} as
an Agent~B proof claim at the exact-adjoint-benchmark level, with independent
re-audit of the matching-reconstruction steps still outstanding.
\end{remark}

% PROV: agent-B/notes/finite-dimensional-adjoint-jb-splitting-corollary.md.
\begin{corollary}[Master adjoint-benchmark corollary]\label{cor:adjoint-benchmark}
\status{(benchmark; modulo \Cref{rem:matrix-audit})}
There is a universal constant $C$ such that for every finite-dimensional
JB-algebra $\mathcal B$ and every exact adjoint coboundary $f=d^1h\in
C^2(\mathcal B,\mathcal B)$ there is a derivation $\delta$ with
$\norm{h-\delta}\le C\norm{f}_{\mathrm{inj}}$: the exact-adjoint coboundary
inversion is dimension-free, uniformly over all $\mathcal B$ and all numbers of
simple summands.
\end{corollary}

\subsection{The next-arrow estimate and what remains open}\label{ssec:layer1-open}

The exact-adjoint splitting (\Cref{cor:adjoint-benchmark}) is the right inverse
for \emph{exact} cocycles. Error reduction additionally needs the
\emph{next-arrow} estimate $\operatorname{dist}(\theta,\Img d^1)\le K'\norm{J\theta}$
for \emph{approximate} cocycles, and several further robustness properties.
Current status of the next-arrow estimate: proved for commutative scalar modules
(\Cref{prop:comm-scalar}), for the normalized spin sector
(\Cref{prop:spin-splitting}), and for the diagonal-source ($D\times D$) block of
matrix factors; \emph{open} for the mixed ($D\times E$) and off-diagonal
($E\times E$) matrix blocks.

% PROV: agent-B/notes/diagonal-frame-matrix-next-arrow-walsh-target.md.
\begin{proposition}[Fixed diagonal-frame matrix next-arrow]\label{prop:diag-next-arrow}
\status{(proved; fixed-frame benchmark)}
Let $D=\R^n$ act on $M=H_n(F)$, $F\in\{\R,\C,\mathbb H\}$, by the diagonal-frame
Jordan action $x\cdot m=x\jp m$. For every unit-normalized symmetric
$2$-cochain $\theta\colon D\times D\to M$,
\[
  \operatorname{dist}(\theta,\Img d^1)\le C\norm{J\theta}_{\mathrm{inj}},
  \qquad
  \norm{\theta-\Pi_n\theta}_{\mathrm{inj}}\le C'\norm{J\theta}_{\mathrm{inj}},
\]
where $\Pi_n=d^1S_n$ is the bounded Rademacher projection onto coboundaries.
The proof decomposes the off-diagonal entries into endpoint, one-tail, and
tail-tail Walsh pieces and reconstructs the tail-tail component by a two-density
sparse-sign identity. This closes the fixed-frame $D\times D$ source block only;
it is not the full matrix next-arrow theorem.
\end{proposition}

For the full high-rank matrix next-arrow problem, a fixed frame splits the source
space as $J=D\oplus E$ (diagonal plus off-diagonal Peirce sectors). The
$D\times D$ block is \Cref{prop:diag-next-arrow}. The open blocks are
$D\times E$, expected to require cochain-level leakage control using diagonal
sign spectral gaps, and $E\times E$, expected to require Peirce
curvature/matching tests.

% PROV: agent-B/notes/layer1-after-adjoint-benchmark-obligations.md, full-matrix-next-arrow-source-decomposition-target.md, frame-covariance-and-global-matrix-obstacle.md, layer1-output-requirements.md.
\begin{openproblem}[What separates the benchmark from \Cref{op:jordan-structure}]
\label{op:layer1-gap}
\status{(open)}
To upgrade \Cref{cor:adjoint-benchmark} to the structure theorem one needs, all
with dimension-free order-unit constants: (i) the pre-cohomological construction
of approximate idempotent elements, Jordan frames, Peirce splittings, and
coordinatization maps from an arbitrary $\eps$-JB algebra; (ii) the next-arrow
estimate for \emph{arbitrary} relevant Jordan modules, not just the adjoint
module (the mixed Peirce-$\tfrac12$ components of \Cref{rem:directsum-scope} and
the matrix $D\times E$, $E\times E$ blocks); (iii) tolerance of
\emph{approximate} module actions, since during the Newton iteration the
codomain is itself only an $\eps$-JB algebra; (iv) a stable gauge across
iterations; (v) very possibly an \emph{incremental} assembly (applying error
reduction only to controlled subalgebras at each merge/extend step), since naive
frame-averaging of the fixed-frame matrix splitting is ruled out --- the
averaged diagonal compression acts on the traceless part by a factor $\sim1/n$,
a rank-sized loss; and (vi) for the \emph{factorization} application, a
positivity-capable or concrete output (\Cref{rem:layer1-route}), which an
abstract Jordan isomorphism does not supply.
\end{openproblem}

In short: the dimension-free order-unit splitting of the Jordan coboundary, long
the suspected obstruction, is now available at the exact-adjoint benchmark level
for every finite-dimensional JB-algebra (\Cref{cor:adjoint-benchmark}); the
structure theorem \Cref{op:jordan-structure} remains open because the
pre-cohomological construction, approximate-cocycle, arbitrary-module,
approximate-module, and positivity-output gaps recorded in \Cref{op:layer1-gap}
are still unresolved. This is genuine and substantial progress over the earlier
state, in which even the exact splitting was unresolved beyond small dimension.
